\chapter{Введение}

\section{Актуальность и постановка проблемы}
Диссертация исследует вычислительные и алгоритмические свойства режимов
предиктивного кодирования в Torch2PC в сопоставлении с обратным
распространением ошибки (BP). Исходное сравнение качества и вычислительной
стоимости последовательно расширяется послойной диагностикой, механизмным
анализом и проверкой возможности безопасного досрочного решения.

Исследование не предполагает превосходство какого-либо режима заранее.
Положительные, отрицательные и смешанные результаты сохраняются как результаты
заранее ограниченных процедур, а границы выводов отделяются от последующих
интерпретаций.

\section{Исследовательские вопросы}
Исследование организовано вокруг трёх вопросов.

\begin{description}
  \item[RQ1.] При каких зафиксированных алгоритмических и вычислительных
  условиях режимы Exact, FixedPred и Strict близки к BP, а при каких различия
  выходят за заранее заданные численные или статистические границы?

  \item[RQ2.] Какие вычислительные механизмы и компоненты стоимости определяют
  наблюдаемые различия между исследованными режимами предиктивного кодирования?

  \item[RQ3.] Можно ли по доступному до завершения вычисления состоянию
  безопасно распознавать достаточные состояния и получить положительную
  вычислительную выгоду от раннего решения?
\end{description}

\section{Логика исследования}
Работа строится как последовательность всё более узких проверок:
Stage~1/2 устанавливают базовые сравнительные результаты; Stage~3A исследует
послойные градиенты и представления; Stage~3B локализует вычислительные
механизмы и стоимость; QWake проверяет, можно ли превратить доступный
pre-action сигнал в безопасное и экономически полезное решение.

\section{Контракт утверждений}
Таблица ниже фиксирует состояние основных утверждений диссертации. Она
генерируется из \texttt{thesis/data/research\_claims.json}; численные QWake C2
границы связаны с отдельным thesis-facing summary, который хранит идентичности
запечатанного результата и не является новым научным evidence.

\input{generated/claims_matrix}

\section{Вклад и границы работы}
Вклад диссертации состоит не только в сравнении режимов, но и в причинной
декомпозиции наблюдаемых различий: от поведения моделей к послойным
механизмам, затем к стоимости наблюдения и к ограниченной проверке безопасного
решения. Утверждения относятся только к явно указанным реализациям,
архитектурам, наборам данных, вычислительной среде и замороженным протоколам.
Отдельно отмечаются вопросы, которые текущим дизайном не проверялись.
